\section{Experimental Methodology}

\subsection{Environment and Controls}

PostgreSQL, Redis, Express, and k6 ran as local processes communicating over loopback. The Linux environment exposed nine logical CPUs and approximately 9.7\,GiB of RAM and reported an AMD EPYC 9V74 processor. The server versions were PostgreSQL 17.11, Redis 8.0.2, Node.js 24.22.1, and Express 4.21.2, with k6 2.3.0 driving the performance runs. Each trial recreated its services and used the same pool limit of ten connections, database schema, stock deduction logic, and dataset loader. Trial manifests record seeds, software versions, and code fingerprints. The Docker Compose configuration is provided for replication.

The trial dataset contains 100 menu items, 150 ingredients, and 300 recipe links. Items are arranged in groups of four sharing two ingredients. Each item also uses a unique ingredient. There are 20 hot items, which receive approximately 80\% of requests in the performance workload. Within each five-iteration block a deterministic shuffle chooses one one-unit order and four item GETs. The dataset loader uses the same seed across paired policies within each repetition. Performance stock starts at one million units per ingredient so that the arrival rate can be varied without depleting inventory. Correctness tests that require a sellout use separate trial datasets.

\subsection{Correctness Exercises}

The induced-refill test compares A, C, and D in ten matched repetitions. It starts with two portions available for an item group and uses the first sale to leave one. The driver starts a GET for a related item on an empty cache key. A loopback-only test gate pauses that GET after its database read. An order sells the last portion and returns its acknowledgment. The driver then releases the delayed GET and issues a new one. The outcome is the version of the new GET compared with authoritative stock after the sale. The experiment records the original delayed response separately. Test hooks are disabled during performance and unscripted freshness runs.

The stockout exercise compares all modes in ten repetitions. It primes the related item with one portion available and sells the last portion in a confirmed competing order. Eight concurrent browsers then read the related item. A browser only attempts an order when its display shows positive stock. A rejection is attributable to a stale display under this schedule when PostgreSQL was already sold out at the start of browsing, the result overstated that stock, and the order returned the insufficient-inventory code. Counts use browsers as the denominator. Orders share inventory within a repetition, so each browser is not a separate experimental replicate.

A separate hot-item exercise uses six reader threads and two order threads for 1.5\,s per repeat, with ten repeats per mode, a three-second TTL, and the high-stock trial dataset. A shared client records the highest acknowledged stock version. Each GET reads that reference before issuing its request. Eligible responses are successful GETs with a known acknowledgment. A lower returned version establishes staleness relative to that acknowledgment. The stress window is shorter than the TTL. In B, a primed payload can therefore survive the whole window. Trial counts and inventory audits are included for all three exercises.

\subsection{Open-Loop Comparison and Uncertainty}

The performance driver uses k6's constant-arrival-rate executor with 128 preallocated virtual users, one request per iteration, and no sleep in the request path \cite{k6arrival}. A two-second warm-up runs 300 item reads per second. A one-trial-per-mode screen measures three-second windows at 200, 600, 1{,}200, 2{,}400, 4{,}800, and 8{,}000 offered requests/s. The reference rate is 600 requests/s. The high-rate rule selects the first screening rate at which any policy delivers under 95\% of offered requests within the window or has overall HTTP P95 above 50\,ms. Dropped iterations indicate that k6 could not start scheduled work with the available VUs. The driver and services share a host, so the screen characterizes the combined test setup \cite{k6drops}.

The comparison uses ten fresh repetitions for each policy and selected rate, measured for five seconds at a 3\,s TTL. The TTL sweep holds arrival rate at 600 requests/s and compares B and D at 1, 3, and 10\,s TTLs, with ten repetitions and 12\,s measurement windows per cell. The schedule rotates the policy order and shuffles cell order within each repetition. Windows are short, and the 10\,s cell observes little more than one TTL interval.

Endpoint P95 is the nearest-rank 95th percentile of k6 send, wait, and receive time among successful requests, including requests that finish during the drain \cite{k6metrics}. Goodput counts successful responses completed inside the window. Delivery divides goodput by the offered rate. The runner separately records dropped iterations, failures, incomplete requests, and the small discrepancy between nominal schedules and starts plus drops near time boundaries. PostgreSQL statement counts are from \texttt{pg\_stat\_statements} \cite{pgstats}. Cache rates use API counters so that D's generation lookups and Lua commands do not enter the payload-hit denominator. Redis INFO snapshots provide server-side diagnostics \cite{redisinfo}.

The analysis reports medians of per-trial summaries with 95\% percentile bootstrap intervals using 4{,}000 resamples of whole trials. Paired differences resample whole repeat blocks. The intervals describe uncertainty within this host and schedule. A functional smoke check briefly shared the host during the first confirmation block, so the CSV supplement also reports medians with that block omitted. Each primary cell retains all ten trials.
